\section{模型评价与推广}

\subsection{模型优点}

\begin{enumerate}
\item \textbf{守恒型有限体积内核}。以控制面通量差组装的离散格式在机器精度内严格守恒
      （残差 $10^{-13}\sim10^{-16}$），且最内控制体内侧面积结构性为零，天然消除柱坐标轴心
      奇点，无需人工正则化。
\item \textbf{Kirchhoff 积分界面平均}。针对 $D$ 跨量级的强非线性，用解析 Kirchhoff 势保证跨界面
      扩散通量守恒，避免了调和平均被最小值锁死导致的非物理结果。
\item \textbf{干基质量坐标化简}。问题四利用干物质守恒 $\rho_s R^2=\sigma$，把移动边界的几何效应
      凝聚为单一时变因子 $1/R(t)^2$，使收缩问题得以沿用固定域内核，并用附件2 数据独立
      检验了纯径向收缩假设。
\item \textbf{多重交叉验证}。Bessel 解析解、Richardson 收敛阶、守恒审计与独立 BDF 求解器
      四重印证，结果不依赖特定实现。
\end{enumerate}

\subsection{模型局限与推广}

模型假设一维轴对称并忽略端面，对短粗物料或非圆柱形状需扩展到二维/三维；基准未显式计入
蒸发潜热，灵敏度分析显示其可使 $t^*$ 延长约 $9\%$，工程应用中可按实测散热补入潜热汇。
收缩按纯径向均匀处理，若物料出现开裂或各向异性收缩，需引入应力—收缩耦合。方法层面，
本文的守恒型有限体积—Kirchhoff—干基坐标框架可直接推广到其他多孔介质干燥、木材与食品
脱水等耦合传热传质问题\upcite{mayor2004,perre_dns}；对边界与物性数据更丰富的场景，还可与
物理信息神经网络等数据—机理融合方法结合，用于快速反演扩散参数与在线预测干燥终点\upcite{si_pinn,malekjani_sdd}。

% DATA_CHECK_PASSED
